% =================================================
% Учительская версия рабочего листа
% =================================================

\worksheettitle
  {Точки на отрезке}
  {Учительская версия: ответы и методические комментарии}

\begin{teacherbox}[Цель листа]
  Научить устанавливать порядок точек, различать целый отрезок и его части,
  находить расстояние между двумя внутренними точками и рассматривать все
  допустимые положения точки, если задача допускает несколько случаев.
\end{teacherbox}

\begin{checkpointbox}[Ответ к входной разминке]
  \[
    CB=AB-AC=11-4=7\text{ см}.
  \]
\end{checkpointbox}

\section{Порядок точек}

\betweenfigure

\begin{propertybox}[Ответ]
  Точка \(C\) лежит между точками \(A\) и \(B\). Порядок:
  \[
    A-C-B.
  \]
\end{propertybox}

\begin{practicebox}[Ответы к заданию 1]
  \begin{enumerate}[label=\alph*),itemsep=0.4em]
    \item Между \(M\) и \(N\) лежит \(K\).
    \item Между \(P\) и \(S\) лежит \(T\).
  \end{enumerate}
\end{practicebox}

\begin{teacherbox}[Методический акцент]
  Просите ученика проговаривать порядок точек до вычислений. Ошибка в
  расположении почти всегда приводит к верному арифметическому действию,
  применённому к неверной модели.
\end{teacherbox}

\section{Целое и части}

\wholepartsfigure

\begin{propertybox}[Главное равенство]
  \[
    AB=AC+CB.
  \]
\end{propertybox}

\begin{practicebox}[Ответы к заданию 2]
  \begin{enumerate}[label=\alph*),itemsep=0.5em]
    \item
      \[
        MN=MK+KN=3\text{ см }8\text{ мм}+5\text{ см }7\text{ мм}
          =9\text{ см }5\text{ мм}.
      \]
    \item
      \[
        MK=MN-KN=14\text{ см}-6\text{ см }5\text{ мм}
          =7\text{ см }5\text{ мм}.
      \]
  \end{enumerate}
\end{practicebox}

\section{Две точки внутри отрезка}

\fourpointsfigure

\begin{propertybox}[Ответ]
  При порядке \(A-C-D-B\):
  \[
    AB=AC+CD+DB,\qquad CD=AB-AC-DB.
  \]
\end{propertybox}

\begin{practicebox}[Ответ к заданию 3]
  \[
    CD=12-4-3=5\text{ см}.
  \]
  Порядок точек: \(A-C-D-B\).
\end{practicebox}

\begin{thinkbox}[Ответ к заданию 4]
  \[
    AC+DB=9+8=17\text{ см}>15\text{ см}=AB.
  \]
  Значит, отрезки, отложенные от разных концов, перекрываются: точка \(D\)
  лежит левее точки \(C\). Расстояние:
  \[
    CD=AC+DB-AB=17-15=2\text{ см}.
  \]
\end{thinkbox}

\begin{teacherbox}[Не вводите формулу раньше модели]
  Формула \(CD=AC+DB-AB\) легко запоминается механически. Сначала полезно
  отметить точки на схеме и увидеть перекрытие двух частей отрезка.
\end{teacherbox}

\section{Расстояния от одного конца}

\sameendpointfigure

\begin{practicebox}[Ответ к заданию 5]
  Ближе к \(M\) расположена точка \(P\). Так как обе длины отложены от
  одной точки \(M\),
  \[
    PQ=MQ-MP=8\text{ см }2\text{ мм}-3\text{ см }5\text{ мм}
      =4\text{ см }7\text{ мм}.
  \]
\end{practicebox}

\begin{teacherbox}[Диагностика]
  Если ученик складывает \(MP\) и \(MQ\), спросите: «Входит ли отрезок
  \(MP\) в отрезок \(MQ\)?» Это возвращает внимание к отношению части и
  целого.
\end{teacherbox}

\clearpage
\section{Два возможных положения}

\twopositionsfigure

\begin{practicebox}[Ответ к заданию 6]
  Точка \(C\) находится на расстоянии \(6\text{ см}\) от \(A\). От неё
  откладываем \(2\text{ см}\) в обе стороны.

  \studenttwocasesanswersfigure

  В первом случае \(AD=4\text{ см}\), поэтому \(BD=6\text{ см}\).
  Во втором случае \(AD=8\text{ см}\), поэтому \(BD=2\text{ см}\).
\end{practicebox}

\begin{checkpointbox}[Ответы к заданию 7]
  \begin{enumerate}[label=\alph*),itemsep=0.45em]
    \item Два положения: \(AD=2\text{ см}\) или \(AD=8\text{ см}\).
    \item Одно положение: откладывание на \(3\text{ см}\) влево выводит
      точку за пределы отрезка, вправо даёт \(AD=4\text{ см}\).
    \item Решений нет: от точки \(C\) невозможно отложить \(11\text{ см}\)
      и остаться на отрезке длиной \(10\text{ см}\).
  \end{enumerate}
\end{checkpointbox}

\candidatefilterfigure

\section{Типичные ошибки}

\begin{center}
  \renewcommand{\arraystretch}{1.42}
  \begin{tabularx}{\textwidth}{
    >{\raggedright\arraybackslash}p{0.30\textwidth}
    >{\raggedright\arraybackslash}p{0.31\textwidth}
    >{\raggedright\arraybackslash}X}
    \toprule
    \rowcolor{TableHeader}
    Ошибка & Что не понято & Наводящий вопрос \\
    \midrule
    Сразу выполняется сложение или вычитание
      & Не установлено расположение точек
      & В каком порядке идут точки? \\
    \addlinespace[0.2em]
    \(AC\) и \(AD\) складываются
      & Не замечено, что обе длины начинаются в \(A\)
      & Входит ли меньший отрезок в больший? \\
    \addlinespace[0.2em]
    Рассмотрен только один вариант точки \(D\)
      & Расстояние можно отложить в двух направлениях
      & Что получится по другую сторону от \(C\)? \\
    \addlinespace[0.2em]
    Сохранён вариант вне \(AB\)
      & Не проверены ограничения задачи
      & Должна ли точка лежать на отрезке? \\
    \bottomrule
  \end{tabularx}
\end{center}

\begin{teacherbox}[Итоговая формулировка ученика]
  Хороший ответ на рефлексию: «Сначала нужно определить порядок точек и
  понять, какие отрезки являются частями каких. Если положение точки не
  задано однозначно, нужно рассмотреть все допустимые случаи».
\end{teacherbox}
